%iffalse
\let\negmedspace\undefined
\let\negthickspace\undefined
\documentclass[journal,12pt,twocolumn]{IEEEtran}
\usepackage{cite}
\usepackage{amsmath,amssymb,amsfonts,amsthm}
\usepackage{algorithmic}
\usepackage{graphicx}
\usepackage{textcomp}
\usepackage{xcolor}
\usepackage{txfonts}
\usepackage{listings}
\usepackage{enumitem}
\usepackage{mathtools}
\usepackage{gensymb}
\usepackage{comment}
\usepackage[breaklinks=true]{hyperref}
\usepackage{tkz-euclide} 
\usepackage{listings}
\usepackage{gvv}                                        
%\def\inputGnumericTable{}                                 
\usepackage[latin1]{inputenc}                                
\usepackage{color}                                            
\usepackage{array}                                            
\usepackage{longtable}                                       
\usepackage{calc}                                             
\usepackage{multirow}                                         
\usepackage{hhline}                                           
\usepackage{ifthen}                                           
\usepackage{lscape}
\usepackage{tabularx}
\usepackage{array}
\usepackage{float}


\newtheorem{theorem}{Theorem}[section]
\newtheorem{problem}{Problem}
\newtheorem{proposition}{Proposition}[section]
\newtheorem{lemma}{Lemma}[section]
\newtheorem{corollary}[theorem]{Corollary}
\newtheorem{example}{Example}[section]
\newtheorem{definition}[problem]{Definition}
\newcommand{\BEQA}{\begin{eqnarray}}
\newcommand{\EEQA}{\end{eqnarray}}
\newcommand{\define}{\stackrel{\triangle}{=}}
\theoremstyle{remark}
\newtheorem{rem}{Remark}

% Marks the beginning of the document
\begin{document}
\bibliographystyle{IEEEtran}
\vspace{3cm}

\title{Chapter 4 Permutations and Combination}
\author{EE24BTECH11012 - Bhavanisankar G S}
\maketitle
\newpage
\bigskip

\renewcommand{\thefigure}{\theenumi}
\renewcommand{\thetable}{\theenumi}

 6 A student is to answer 10 out of 13 questions in an examination such that he must choose atleast four from the first five questions. The number of choices available to him is  \hfill[2003]
\begin{enumerate} 
    \item  346
    \item  140 
    \item  196 
    \item  280
\end{enumerate}
7.The number of ways in which 6 men and 5 women can dine around a round table if no two women are to sit together is given by \hfill[2003]
\begin{enumerate}
 \item $7!\times5!$
 \item $6!\times 5!$
 \item $30!$
 \item $5!\times4!$
 \end{enumerate}
 8. How many ways are there to arrange the letters in the word GARDEN with vowels iin alphabetical order \hfill[2004]
 \begin{enumerate}
     \item 480
     \item 240
     \item 360
     \item 120
 \end{enumerate}
 9. The number of ways of distributing 8 identical balls in 3 distinct boxes so that none of the boxes is empty is \hfill[2004]
 \begin{enumerate}
     \item $ 8\choose3$
     \item 21
     \item $3^8$
     \item 5
 \end{enumerate}
 10. If the letters of the word SACHIN are arranged in all the possible ways and these words are written out as in dictionary, then the word SACHIN appears at the serial number \hfill[2005]
 \begin{enumerate}
     \item 601
     \item 600
     \item 603
     \item 602
 \end{enumerate}
11. At an election, a voter may vote for any number of candidates, not greater than the number to be elected. There are 10 candidates and 4 are of be selected, if a voter votes for atleast one candidate, then the number of ways in which he can vote is \hfill[2006]
\begin{enumerate}
 \item 5040
 \item 6210
 \item 385
 \item 1110
 \end{enumerate}
 12. The set $S={1,2,3,...,12}$ is to be partitioned into three sets A,B and C of equal size. Thus, $A\bigcup B\bigcup C=S$, $A \bigcap B=B \bigcap  C =A \bigcap C=\phi$. The number of ways to partition S is \hfill[2007]
 \begin{enumerate}
     \item $ \frac{12!}{(4!)^3}$
     \item $\frac{12!}{(4!)^4}$
     \item $\frac{12!}{3!(4!)^3}$
     \item $\frac{12!}{3!(4!)^4}$
 \end{enumerate}
 13. In a shop, there are five types of ice-creams available.A child buys six ice-creams.\\\textbf{Statement1}: The number of different ways in which the child can buy six ice-creams is $10\choose5$.\\\textbf{Statement2}:The number of different ways in which the child can buy six ice-creams is equal to the number of different ways of arranging 6 A's and 4 B's in a row. \hfill[2008]
 \begin{enumerate}
     \item Statement 1 is false, Statement 2 is true.
     \item Statement 1 is true, Statement 2 is true, Statement 2 is the correct explanation of Statement 1.
     \item Statement 1 is true, Statement 2 is true, Statement 2 is not a correct explanation of statement 1.
     \item Statement 1 is true, Statement 2 is false.
 \end{enumerate}
 14.How many different words can be formed by jumbling the letters in the word MISSISSIPPI in which no two S are adjacent ?
 \begin{enumerate}
     \item 8$6\choose4$$7\choose4$
     \item $6\times7$$8\choose4$
     \item $6\times8$$7\choose4$
     \item 7$6\choose4$$8\choose4$
 \end{enumerate}
 15. From 6 different novels and 3 different dictionaries, 4 novels and 1 dictionary are to be selected and arranged in a row on a shelf so that the dictionary is always in the middle.Then the number of such arrangements is: \hfill[2009]
 \begin{enumerate}
     \item at least 500 but less than 750
     \item at least 750 but less than 1000
     \item at least 1000
     \item less than 500
 \end{enumerate}
 16. There are two urns.Urn A has 3 distinct red balls while Urn B has 9 distinct blue balls.From each urn, two balls are taken at random and then transferred to the other.The number of ways in which this can be done is: \hfill[2010]
 \begin{enumerate}
     \item 36
     \item 66
     \item 108
     \item 3
 \end{enumerate}
 17.\textbf{Statement 1}: The number of ways of distributing 10 identical balls in 4 identical boxes such that no box is empty is $9\choose3$.\\\textbf{Statement 2}:The number of ways of choosing any three 3 places from 9 different places is $9\choose3$. \hfill[2011]
 \begin{enumerate}
     \item Statement 1 is true, statement 2 is true; Statement 2 is the correct explanation of Statement 1.
     \item Statement 1 is true, Statement 2 is true; Statement 2 is not the correct explanation of Statement 1.
     \item Statement 1 is true, Statement 2 is false.
     \item Statement 1 is false, Statement 2 is true.
 \end{enumerate}
 18.These are 10 points in a plane, out of which 6 are collinear.If N is the number of triangles formed by these points, then; \hfill[2012]
 \begin{enumerate}
     \item $N\le100$
     \item $100<N\le140$
     \item $140<N\le190$
     \item $N>190$
 \end{enumerate}
 19. Assuming the balls to be identical except for the difference in colours, The number of ways in which one or more balls can be selected from 10 white, 9 green and 7 black balls. \hfill[2012]
 \begin{enumerate}
     \item 880
     \item 629
     \item 630
     \item 879
 \end{enumerate}
 20. Let $T_n$ be the set of all possible triangles formed by joining the vertices of a n-sided regular polygon. If $T_{n+1} - T_n = 10$, then the value of n is : \hfill[JEE M 2013]
 \begin{enumerate}
     \item 7
     \item 5
     \item 10
     \item 8
 \end{enumerate}
\end{document}
